\documentclass[conference]{IEEEtran}
\usepackage[T1]{fontenc}
\usepackage[utf8]{inputenc}
\usepackage[brazil]{babel}
\AtBeginDocument{\shorthandoff{"}}
\usepackage{amsmath,amssymb}
\usepackage{booktabs,tabularx,array}
\usepackage{cite}
\usepackage{url}
\usepackage{tikz}
\usetikzlibrary{positioning,arrows.meta}
\newcolumntype{L}{>{\raggedright\arraybackslash}X}
\newcolumntype{P}[1]{>{\raggedright\arraybackslash}p{#1}}
\newcommand{\cod}[1]{{\ttfamily\def\_{\textunderscore\allowbreak}#1}}
\renewcommand{\arraystretch}{1.08}

\begin{document}
\title{UX-Agent: Implementação de um Agente Proxy Baseado em LLM com Paridade de Signos, Mecanismos Fora do Prompt e Verificação Externa para um Chat de Estudo}

\author{
\IEEEauthorblockN{Matheus Emanuel da Silva}
\IEEEauthorblockA{Engenharia de Computação\\CEFET MG, Campus Divinópolis\\matheus.silva@aluno.cefetmg.br}
\and
\IEEEauthorblockN{João Francisco Teles da Silva}
\IEEEauthorblockA{Engenharia de Computação\\CEFET MG, Campus Divinópolis\\joao.silva.05@aluno.cefetmg.br}
\and
\IEEEauthorblockN{Eduardo Almeida}
\IEEEauthorblockA{Engenharia de Computação\\CEFET MG, Campus Divinópolis\\eduardo.almeida@aluno.cefetmg.br}
}
\maketitle

\begin{abstract}
Agentes baseados em grandes modelos de linguagem que simulam participantes costumam encenar o perfil pelo prompt, recebem informação que a pessoa não teria e declaram o próprio sucesso. Este artigo relata a implementação do UX-Agent para um chat de estudo usado por uma persona de estudante universitária autista e mostra que decisões de design tornam os registros auditáveis. A tela humana e a observação textual do agente são geradas pela mesma função, garantindo paridade de signos; o diagrama MoLIC é a máquina de estados e restringe a resposta do modelo, por esquema JSON, às falas do usuário habilitadas em cada cena; cinco mecanismos de perfil, a hesitação e o limite de tentativas são sorteados com semente fora do prompt; a ajuda, em três estratos, só chega por ação registrada; e um verificador externo decide o sucesso pelo estado final. Rupturas de comunicação e mecanismos de prevenção e recuperação são detectados apenas por assinaturas no log. No piloto, com 120 sessões das tarefas de calibração e o modelo gpt-4o-mini, não houve proposta inválida nem falha técnica; a tarefa de localização teve sucesso total nas duas condições, e a tarefa com metáfora e ironia teve sucesso nulo também no comparador, o que revelou uma chave ambígua. O piloto também mostrou um mecanismo de suporte inoperante e a ausência de dispersão no comparador. O tempo simulado ficou de 22\% a 39\% acima da linha de base KLM. As tarefas de validação permanecem intocadas.
\end{abstract}

\begin{IEEEkeywords}
Interação humano-computador, agente baseado em LLM, MoLIC, paridade de signos, KLM.
\end{IEEEkeywords}

\section{Introdução}
Agentes baseados em grandes modelos de linguagem (LLM) já são usados para simular participantes em testes de usabilidade \cite{lu2025uxagent}. Três falhas recorrentes tornam esses registros pouco confiáveis: o perfil é apenas encenado por instruções no prompt, o agente recebe informação que a pessoa real não teria e o sucesso é declarado pelo próprio modelo \cite{lara2026p3}. Estudos de fidelidade mostram que LLMs instruídos por identidade homogeneízam e estereotipam grupos \cite{wang2025,park2025} e reproduzem mal ações individuais \cite{lu2026}.

Este artigo relata a implementação do UX-Agent para um chat de estudo voltado a estudantes universitários autistas e responde à pergunta-chave: \emph{que decisões de design da interface instrumentada, do sistema de ajuda e do acoplamento agente-sistema permitem gerar registros auditáveis fornecendo ao agente apenas a informação que a persona teria?} A resposta é uma arquitetura em que (i) tela humana e observação do agente saem da mesma função, (ii) o diagrama MoLIC é a máquina de estados e restringe a resposta do LLM, (iii) o comportamento da persona vem de mecanismos parametrizados no código, (iv) a ajuda só chega por ação registrada e (v) o sucesso é decidido por um verificador externo. O caderno de design (Anexo E1) especifica metamensagem, paridade de signos, CTT, KLM, MoLIC, matriz de prevenção e recuperação e ajuda. O piloto usa só as tarefas de calibração.

\section{Modelagem Herdada}
A persona E. é fictícia: 21 anos, Engenharia de Computação, estuda sozinha à noite, adulta autista sem deficiência intelectual e com linguagem funcional (CID-11 6A02.0) \cite{who2024}, usuária regular de web e chat. O comparador sintético é o mesmo agente com os mecanismos de perfil desligados. Quatro tarefas foram congeladas na modelagem: C1 (localizar local e horário numa fala curta) e C2 (metáfora e ironia numa fala longa com bloco oculto) para calibração; V1 (localizar sala e prazo entre oito distratores) e V2 (instrução implícita) para validação. Cada tarefa tem uma HTA \cite{annett1967} com objetivo, planos e operações ligadas a eventos (ler cartão, rolar, reformular, nova pergunta, enviar campo, encerrar), caminho mínimo $A_{\min}$ (3 em C1; 4 nas demais) e erros previstos.

As direções dos mecanismos vêm da literatura: a diferença de compreensão figurada é maior na metáfora que na ironia e quase some quando a habilidade de linguagem é equiparada \cite{kalandadze2018}; em busca na web, adultos autistas exploram mais elementos com sucesso semelhante \cite{eraslan2019}. As magnitudes são hipóteses operacionais.

\section{Decisões de Design}
A Tabela~\ref{tab:dec} lista as decisões com justificativa e status epistemológico: E (evidência empírica), P (princípio de design) ou H (hipótese operacional).

\begin{table*}[t]
\caption{Decisões de design, justificativa e status (E: evidência; P: princípio de design; H: hipótese operacional).}
\label{tab:dec}
\centering\footnotesize
\begin{tabularx}{\textwidth}{@{}P{0.5cm}P{5.3cm}LP{0.4cm}@{}}
\toprule
Id & Decisão & Justificativa & St.\\
\midrule
D1 & Tela humana e observação $o_t$ geradas pela mesma função \cod{tela()} & Paridade de signos: a interface é metacomunicação com dois receptores \cite{kammersgaard1988,desouza2009} & P\\
D2 & Resposta do LLM forçada por esquema JSON cujo \cod{enum} são as falas $u{:}$ habilitadas na cena MoLIC & A MoLIC define o espaço de ações; o agente não consegue agir fora do diálogo projetado \cite{barbosa2003} & P\\
D3 & HTA, CTT, KLM e gabarito no pacote \cod{analista/}, nunca importado pelo agente & Fronteira de informação: instrumentos do analista \cite{dix2004,mackenzie2013} & P\\
D4 & Filtro de leitura (R1) modulado pelo tipo de figura e pela habilidade de linguagem (R5) & Direção: metáfora $>$ ironia; efeito some com linguagem equiparada \cite{kalandadze2018}; magnitude H & E/H\\
D5 & Injetor de exploração (R2): retira enviar e encerrar por um passo antes do primeiro envio & Mais exploração com sucesso semelhante \cite{eraslan2019}; magnitude H & E/H\\
D6 & Limiar de suporte (R3) e aceite prematuro (R4) como filtros de \emph{affordance} & Desejo de esclarecimento sob demanda e dificuldade de avaliar respostas relatados em estudos de percepção \cite{haroon2024,jang2024} & H\\
D7 & Hesitação lognormal por passo, +1 $M$ na primeira exposição a figura; desistência após $k$ ações sem progresso & Requisito do enunciado (tempo de hesitação, limite de tentativas) \cite{laviola2017}; sem estimativa publicada & H\\
D8 & Ajuda em três estratos (significado, procedimento, recuperação), só por fala registrada & Ajuda sob demanda por tipo de dúvida \cite{shneiderman2010} & P\\
D9 & Confirmação ao encerrar com campo vazio (PA), recusa de texto vazio (CE), substituição com valor anterior (RA), alerta com 3 ações (AL) & Prevenção e recuperação de erros \cite{shneiderman2010,reason1990} & P\\
D10 & Rupturas MAC e mecanismos PP/PA/AL/RA/CE detectados só por assinatura de eventos & Não inferir ruptura do discurso emocional do LLM \cite{lara2026p3,desouza2009,nielsen2010} & P\\
D11 & Sucesso decidido por verificador que lê só o estado final & A afirmação do agente não é critério \cite{lara2026p3} & P\\
D12 & Tempo simulado = operadores KLM da fala + hesitação; latência da API gravada à parte & Operadores de \cite{card1980}; hesitação H & E/H\\
D13 & V1 e V2 bloqueadas no código sem a opção \cod{--validacao-congelada} & Separação calibração/validação & P\\
D14 & Temperatura 0 e semente fixa; memória das próprias ações dentro da tarefa, reiniciada entre tarefas & Reprodutibilidade; memória de curto prazo equivalente à da pessoa & P/H\\
\bottomrule
\end{tabularx}
\end{table*}

\section{Arquitetura e Acoplamento Agente-Sistema}
A Fig.~\ref{fig:arq} mostra o ciclo de um passo. O \emph{amostrador} sorteia, no início da execução, $h$, $p_{\text{lit}}$, $p_{\text{explore}}$, $p_{\text{sup}}$, $p_{\text{accept}}$, a hesitação média e o limite $k$ (semente $=$ 20261026 $+\,i$). A cada passo: (1) a MoLIC devolve as falas habilitadas no estado; (2) o controlador aplica os filtros de \emph{affordance}; (3) o tradutor monta $o_t$ a partir da tela, aplicando o filtro de leitura se sorteado; (4) o LLM (\cod{gpt-4o-mini-2024-07-18}, temperatura 0) recebe prompt de sistema fixo, instrução do protocolo, histórico das próprias ações e $o_t$, e responde \cod{\{acao, campo, texto\}} sob esquema estrito; (5) o executor aplica a fala, gera a fala $d{:}$ e atualiza o estado; (6) o registrador grava os eventos. Propostas fora da lista são rejeitadas; três seguidas, ou três erros de API, encerram como falha técnica. O verificador roda depois, fora do laço. O prompt de sistema não menciona autismo, diagnóstico, figura nem critério de sucesso.

\begin{figure}[t]
\centering
\begin{tikzpicture}[font=\scriptsize, b/.style={draw, rounded corners=1pt, align=center, minimum height=7mm, minimum width=17mm, inner sep=2pt}, >=Stealth]
\node[b] (amo) at (0,0) {Amostrador\\$\theta$, semente};
\node[b] (mol) at (2.8,0) {MoLIC\\falas $u{:}$};
\node[b] (ctl) at (5.6,0) {Controlador\\R1--R5, $k$};
\node[b] (obs) at (5.6,-1.3) {Tradutor $o_t$};
\node[b] (llm) at (5.6,-2.6) {LLM\\JSON + \cod{enum}};
\node[b] (exe) at (2.8,-2.6) {Executor\\$u{:}\rightarrow d{:}$};
\node[b] (est) at (2.8,-1.3) {Estado da\\interface};
\node[b] (log) at (0,-2.6) {Registrador\\JSON};
\node[b] (ver) at (0,-1.3) {Verificador\\externo};
\draw[->] (amo.north) -- ++(0,0.35) -| (ctl.north);
\draw[->] (mol) -- (ctl);
\draw[->] (ctl) -- (obs);
\draw[->] (est) -- (obs);
\draw[->] (obs) -- (llm);
\draw[->] (llm) -- node[above]{proposta} (exe);
\draw[->] (exe) -- (est);
\draw[->] (est) -- (mol);
\draw[->] (exe) -- (log);
\draw[->, dashed] (est) -- node[above]{final} (ver);
\draw[->, dashed] (ver) -- (log);
\end{tikzpicture}
\caption{Acoplamento agente-sistema. Tracejado: fora do laço do agente.}
\label{fig:arq}
\end{figure}

\section{Instrumentação}
Cada execução gera um arquivo JSON com \cod{meta}, \cod{eventos} e \cod{estado\_final}. Todo evento traz identificador da execução, persona, tarefa, conjunto, condição, \cod{config\_id} (hash de $\theta$, prompt, roteiro e modelo), semente, passo, carimbo de tempo ISO 8601, tipo, cena MoLIC, fala $u{:}$, resposta $d{:}$, dados e origem. São 15 tipos, entre eles \cod{observacao\_exposta} (tela integral), \cod{observacao\_transformada} ($o_t$ enviado), \cod{restricao\_aplicada} (mecanismo, falas retiradas e probabilidade), \cod{acao\_proposta} (prompt completo, resposta bruta, modelo devolvido, tokens e latência), \cod{acao\_executada} (campos antes e depois, operadores KLM, hesitação) e \cod{avaliacao\_externa} (escrito só pelo verificador). Assim, cada desvio entre o que a tela mostrou e o que o agente recebeu fica no próprio log.

Regras de detecção, todas sobre eventos: deslize = \cod{erro\_capturado} ou campo trocado; lapso = campo obrigatório vazio no fim; engano de interpretação = envio fora da chave num campo ligado a figura; engano de avaliação = \cod{encerrada} com campo preenchido fora da chave \cite{reason1990}. Rupturas MAC \cite{desouza2009}: ``Cadê?'' (enviar antes de rolar com bloco oculto), ``O que é isto?'' (ajuda de significado), ``Epa!'' (campo substituído), ``Assim não dá'' (nova pergunta após reformular), ``E agora?'' (duas ações sem progresso), ``Para mim está bom'' (engano de avaliação) e ``Desisto'' (limite de tentativas, confirmação com campo vazio ou interrupção).

\section{Piloto}
\subsection{Execução}
Foram 120 sessões em 22/11/2026: C1 e C2, condições perfil (mecanismos ligados) e comparador (desligados), 30 execuções cada, sementes 20261026 a 20261055, modelo devolvido pela API \cod{gpt-4o-mini-2024-07-18}. Não houve proposta rejeitada nem falha técnica. Foram 282.403 tokens no total, com latência mediana de 0,74 a 0,91\,s por passo, registrada à parte do tempo simulado. V1 e V2 não foram abertas. Dois executores de regra (resposta canônica e resposta errada) confirmaram antes o verificador: 2/2 e 0/2 sucessos.

\subsection{Métricas}
A Tabela~\ref{tab:res} resume as métricas por tarefa e condição. $S$ é a taxa de sucesso dada pelo verificador; $R_{e,\text{interp}}$ é a fração de oportunidades de figura com engano de interpretação; $R_{e,\text{aval}}$ é a fração de sessões encerradas com campo fora da chave; $A$ é o número de ações executadas; $U_s$ é a fração de sessões com ajuda de significado.

\begin{table}[t]
\caption{Resultados do piloto (n = 30 por célula).}
\label{tab:res}
\centering\footnotesize
\begin{tabularx}{\columnwidth}{@{}lLLLL@{}}
\toprule
 & C1 comp. & C1 perfil & C2 comp. & C2 perfil\\
\midrule
$S$ & 1,00 & 1,00 & 0,00 & 0,00\\
$R_{e,\text{interp}}$ & -- & -- & 0,50 & 0,57\\
$R_{e,\text{aval}}$ & 0,00 & 0,00 & 1,00 & 1,00\\
$A$ mediana (faixa) & 3 (3--3) & 3 (3--4) & 4 (4--4) & 4 (4--5)\\
$A_{\min}$ (HTA/KLM) & 3 & 3 & 4 & 4\\
$U_s$ (sessões) & 0,00 & 0,47 & 0,00 & 0,00\\
$T_{\text{sim}}$ mediano (s) & 22,8 & 24,8 & 33,3 & 34,8\\
$T_{\text{KLM}}$ (s) & 17,8 & 17,8 & 27,4 & 27,4\\
$T_{\text{sim}}/T_{\text{KLM}}$ & 1,28 & 1,39 & 1,22 & 1,27\\
\bottomrule
\end{tabularx}
\end{table}

\emph{C1.} As duas condições tiveram sucesso em todas as sessões. No comparador as 30 sessões seguiram exatamente o caminho mínimo da HTA (enviar, enviar, encerrar). No perfil, o injetor de exploração disparou em 14 sessões e, nas 14, o agente usou \cod{reformular} antes do primeiro envio (assinatura ``O que é isto?''), elevando $A$ para 4. O aceite prematuro disparou em 2 sessões sem efeito no resultado.

\emph{C2.} Nenhuma sessão teve sucesso. No comparador, as 30 sessões preencheram Prazo corretamente (``sexta, 18h'') e Entrega com ``rascunho do texto individual'', fora da chave (``ideia central'' e ``duas consequências''), e encerraram: 30 enganos de avaliação (``Para mim está bom''). No perfil, o filtro de leitura disparou na ironia em 9 sessões e, nas 9, Prazo recebeu ``segunda'', como previsto por R1. Nas 5 dessas sessões em que só a ironia foi filtrada, Entrega ficou correta: sem a frase ``a secretaria confirmou o prazo do rascunho'', o agente deixou de ancorar Entrega em ``rascunho''. O limiar de suporte disparou em 18 sessões, mas o agente escolheu \cod{rolar\_historico} em todas, e o botão de explicação nunca foi usado.

\emph{Tempo.} O tempo simulado ficou de 22\% a 39\% acima do KLM. O excesso vem da hesitação sorteada e, no perfil, das ações extras; o KLM não inclui leitura.

\subsection{O que o piloto depurou}
O piloto cumpriu sua função de depurar o protocolo antes da validação \cite{lara2026p3}. Quatro achados mudam a próxima rodada de calibração.
(i) \emph{Chave de C2 ambígua.} ``Texto individual'' é uma resposta defensável a ``o que deve ser entregue'', e o comparador fica em $S=0$; com piso no comparador, C2 não discrimina condições. A chave ou o cartão de C2 precisa ser revista, o que é permitido por ser tarefa de calibração; V1 e V2 continuam congeladas.
(ii) \emph{R3 inoperante.} Retirar \cod{enviar} e \cod{encerrar} não leva ao suporte quando \cod{rolar\_historico} está disponível; a regra deve atuar só depois de o bloco oculto ser revelado.
(iii) \emph{Confusão entre R2 e ajuda.} Em C1, todo uso de ajuda veio do injetor de exploração; $U_s$ deve ser relatado separando os passos com \cod{restricao\_aplicada}.
(iv) \emph{Comparador sem dispersão.} Com temperatura 0 e observação idêntica, as 30 sessões do comparador foram iguais; a variabilidade só vem dos mecanismos. Isso confirma o risco de homogeneização apontado por \cite{wang2025} e indica que a variação deve ser modelada explicitamente.
Por fim, todas as métricas puderam ser calculadas a partir do log, com denominador definido e sem sessões incompletas.

\section{Ameaças à Validade}
\emph{Vazamentos residuais.} A paridade é por construção, mas não total. (i) A lista de ações traz o id e uma descrição curta de cada botão; a pessoa vê só o rótulo. A descrição é metalinguística e não cita a tarefa, mas é mais explícita que o botão. (ii) A instrução do cartão diz quais campos preencher, o que a pessoa também lê. (iii) A memória do agente é o resumo literal das próprias ações e respostas $d{:}$, mais fiel que a memória humana. (iv) O nome do campo (por exemplo, ``Prazo'') orienta a busca para os dois receptores igualmente.

\emph{Mecanismos.} As direções de R1, R2 e R5 têm fonte; as magnitudes, a hesitação e o limite de tentativas são hipóteses operacionais sem estimativa publicada. Os filtros de \emph{affordance} retiram botões por um passo para o agente, enquanto a pessoa vê todos; esse desvio é intencional e registrado, mas não equivale a um modelo cognitivo.

\emph{LLM.} Com temperatura 0 e mesma observação, o modelo responde de forma quase determinística; a variação entre execuções vem quase toda dos sorteios do controlador, o que concentra a dispersão no comparador. O \cod{seed} da API é de melhor esforço e o \cod{system\_fingerprint} pode mudar. O modelo pode ter visto textos parecidos no treino.

\emph{KLM e tempo.} O KLM não tem operador de leitura, então a linha de base subestima tarefas com texto longo; a regra de um $M$ por fala é uma heurística. O tempo simulado não é tempo humano e não deve ser comparado com estudos com pessoas.

\emph{Escopo.} O piloto usa só C1 e C2 e não estima desempenho humano. V1 e V2 continuam intocadas para a validação.

\section{Conclusão}
A pergunta-chave foi respondida por cinco decisões verificáveis no código e no log: a mesma função gera a tela humana e a observação do agente (paridade de signos); a MoLIC restringe a resposta do LLM às falas $u{:}$ do estado; o comportamento da persona sai de mecanismos sorteados com semente, fora do prompt, e cada intervenção é registrada; a ajuda chega só por fala registrada e estratificada por tipo de dúvida; e o sucesso é decidido por um verificador que lê apenas o estado final. Essas escolhas tornaram auditável cada diferença entre o que a tela mostrou e o que o agente recebeu, e permitiram ao piloto revelar defeitos de protocolo (chave ambígua em C2, R3 inoperante, confusão entre exploração e ajuda) antes de qualquer uso das tarefas de validação ou de estudo com pessoas. Os próximos passos são revisar C2 e R3 na calibração, congelar $\theta$, prompt e modelo, e só então executar V1 e V2. O agente continua um instrumento de depuração e não substitui participantes \cite{lu2025uxagent}.

\begin{thebibliography}{99}
\footnotesize
\bibitem{lu2025uxagent} Y. Lu \emph{et al.}, ``UXAgent: An LLM agent-based usability testing framework for web design,'' in \emph{Ext. Abstr. CHI Conf. Human Factors Comput. Syst.}, Yokohama, 2025, doi: 10.1145/3706599.3719729.
\bibitem{lara2026p3} A. P. Lara, ``Desenvolvimento e implementação do UX-Agent: Projeto de Interação e Implementação do Sistema Instrumentado,'' enunciado da disciplina de IHC, CEFET-MG, Divinópolis, nov. 2026.
\bibitem{wang2025} A. Wang, J. Morgenstern, and J. P. Dickerson, ``Large language models that replace human participants can harmfully misportray and flatten identity groups,'' arXiv:2402.01908v3, 2025.
\bibitem{park2025} S. Park, A. Min, J. A. Beltran, and G. R. Hayes, ``As an autistic person myself: The bias paradox around autism in LLMs,'' in \emph{Proc. CHI}, 2025, doi: 10.1145/3706598.3713420.
\bibitem{lu2026} Y. Lu \emph{et al.}, ``Can LLM agents simulate multi-turn human behavior? Evidence from real online customer behavior data,'' in \emph{Proc. 64th Annu. Meeting ACL}, 2026, pp. 43961--43977.
\bibitem{who2024} World Health Organization, ``6A02 Autism spectrum disorder,'' \emph{ICD-11 for Mortality and Morbidity Statistics}, 2024. [Online]. Disponível: https://icd.who.int/browse/2024-01/mms/en\#437815624.
\bibitem{annett1967} J. Annett and K. D. Duncan, ``Task analysis and training design,'' \emph{Occupational Psychology}, vol. 41, pp. 211--221, 1967.
\bibitem{kalandadze2018} T. Kalandadze, C. Norbury, T. Nærland, and K.-A. B. Næss, ``Figurative language comprehension in individuals with autism spectrum disorder: A meta-analytic review,'' \emph{Autism}, vol. 22, no. 2, pp. 99--117, 2018.
\bibitem{eraslan2019} S. Eraslan, V. Yaneva, Y. Yesilada, and S. Harper, ``Web users with autism: Eye tracking evidence for differences,'' \emph{Behav. Inf. Technol.}, vol. 38, no. 7, pp. 678--700, 2019.
\bibitem{kammersgaard1988} J. Kammersgaard, ``Four different perspectives on human-computer interaction,'' \emph{Int. J. Man-Mach. Stud.}, vol. 28, no. 4, pp. 343--362, 1988.
\bibitem{desouza2009} C. S. de Souza and C. F. Leitão, \emph{Semiotic Engineering Methods for Scientific Research in HCI}. San Rafael, CA: Morgan \& Claypool, 2009.
\bibitem{barbosa2003} S. D. J. Barbosa and M. G. de Paula, ``Designing and evaluating interaction as conversation: A modeling language based on semiotic engineering,'' in \emph{DSV-IS 2003}, LNCS 2844. Berlin: Springer, 2003, pp. 16--33.
\bibitem{dix2004} A. Dix, J. Finlay, G. D. Abowd, and R. Beale, \emph{Human-Computer Interaction}, 3rd ed. Harlow: Pearson, 2004.
\bibitem{mackenzie2013} I. S. MacKenzie, \emph{Human-Computer Interaction: An Empirical Research Perspective}. Waltham, MA: Morgan Kaufmann, 2013.
\bibitem{haroon2024} R. Haroon and F. Dogar, ``TwIPS: A large language model powered texting application to simplify conversational nuances for autistic users,'' in \emph{Proc. ACM ASSETS}, 2024, doi: 10.1145/3663548.3675633.
\bibitem{jang2024} J. Jang, S. Moharana, P. Carrington, and A. Begel, ``It's the only thing I can trust: Envisioning large language model use by autistic workers for communication assistance,'' in \emph{Proc. CHI}, 2024, doi: 10.1145/3613904.3642894.
\bibitem{laviola2017} J. J. LaViola, E. Kruijff, R. P. McMahan, D. A. Bowman, and I. Poupyrev, \emph{3D User Interfaces: Theory and Practice}, 2nd ed. Boston: Addison-Wesley, 2017.
\bibitem{shneiderman2010} B. Shneiderman and C. Plaisant, \emph{Designing the User Interface: Strategies for Effective Human-Computer Interaction}, 5th ed. Boston: Addison-Wesley, 2010.
\bibitem{reason1990} J. Reason, \emph{Human Error}. Cambridge: Cambridge Univ. Press, 1990.
\bibitem{card1980} S. K. Card, T. P. Moran, and A. Newell, ``The keystroke-level model for user performance time with interactive systems,'' \emph{Commun. ACM}, vol. 23, no. 7, pp. 396--410, 1980.
\bibitem{nielsen2010} J. Nielsen and K. Pernice, \emph{Eyetracking Web Usability}. Berkeley, CA: New Riders, 2010.
\end{thebibliography}


\end{document}
